% Môn: Vật lý
% Lớp: 11
% Chương: Chương I. Dao động
% Bài: L11C1 Bài 1. Dao động điều hòa · Xác định các đại lượng đặc trưng và viết phương trình từ đồ thị li độ - thời gian
% ===== Câu 62 =====
% ID: L11C1B1-01-TLN
% Mức: TH
\dangbt{Xác định các đại lượng đặc trưng và viết phương trình từ đồ thị li độ - thời gian}
\begin{ex}
% ID: L11C1B1-01-TLN
% Mức: TH
Trên đồ thị li độ–thời gian, hai đỉnh dương liên tiếp ở $t=1$ s và $t=2$ s. Chu kì theo đơn vị s bằng bao nhiêu?
\shortans{1}
\end{ex}

% ===== Câu 66 =====
% ID: L11C1B1-02-TLN
% Mức: VD
\begin{ex}
% ID: L11C1B1-02-TLN
% Mức: VD
Trên đồ thị li độ–thời gian, hai đỉnh dương liên tiếp ở $t=1$ s và $t=3$ s. Chu kì theo đơn vị s bằng bao nhiêu?
\shortans{2}
\end{ex}

% ===== Câu 70 =====
% ID: L11C1B1-03-TLN
% Mức: VD
\begin{ex}
% ID: L11C1B1-03-TLN
% Mức: VD
Trên đồ thị li độ–thời gian, hai đỉnh dương liên tiếp ở $t=1$ s và $t=5$ s. Chu kì theo đơn vị s bằng bao nhiêu?
\shortans{4}
\end{ex}

% ===== Câu 73 =====
% ID: L11C1B1-07-DS
% Mức: TH
\begin{ex}
% ID: L11C1B1-07-DS
% Mức: TH
Đồ thị li độ–thời gian cho biết $A=4$ cm, hai đỉnh dương liên tiếp cách nhau $1\,\text{s}$ và tại $t=0$ vật ở biên dương. Xác định $T,f,\omega$ và viết phương trình dao động.
\choiceTF

\end{ex}

% ===== Câu 77 =====
% ID: L11C1B1-16-DS
% Mức: VD
\begin{ex}
% ID: L11C1B1-16-DS
% Mức: VD
Đồ thị li độ–thời gian cho biết $A=6$ cm, hai đỉnh dương liên tiếp cách nhau $2\,\text{s}$ và tại $t=0$ vật ở biên dương. Xác định $T,f,\omega$ và viết phương trình dao động.
\choiceTF

\end{ex}

% ===== Câu 81 =====
% ID: L11C1B1-22-DS
% Mức: VD
\begin{ex}
% ID: L11C1B1-22-DS
% Mức: VD
Đồ thị li độ–thời gian cho biết $A=8$ cm, hai đỉnh dương liên tiếp cách nhau $4\,\text{s}$ và tại $t=0$ vật ở biên dương. Xác định $T,f,\omega$ và viết phương trình dao động.
\choiceTF

\end{ex}

% ===== Câu 129 =====
% ID: L11C1B1-23-DS
% Mức: TH
\begin{ex}
% ID: L11C1B1-23-DS
% Mức: TH
Khi phân tích các đặc điểm của đồ thị li độ – thời gian của một vật dao động điều hòa quanh vị trí cân bằng, hãy nhận xét các mệnh đề sau:
\choiceTF
{\True Tọa độ li độ $x$ của vật tại một thời điểm được xác định trực tiếp bằng tung độ của điểm biểu diễn tương ứng trên đồ thị.}
{\True Độ dốc của đồ thị tại một điểm bất kì biểu thị giá trị vận tốc tức thời của vật tại thời điểm đó.}
{\True Nếu đường hình sin của đồ thị có xu hướng đi xuống cắt trục hoành thì vật đang chuyển động qua vị trí cân bằng theo chiều âm.}
{Khoảng cách giữa hai điểm cắt trục hoành liên tiếp của đồ thị luôn có giá trị bằng một chu kì dao động $T$.}
\loigiai{
\begin{itemchoice}
\item (Đúng) Tọa độ li độ $x$ của vật tại một thời điểm được xác định trực tiếp bằng tung độ của điểm biểu diễn tương ứng trên đồ thị. (Vì): Mệnh đề thứ nhất đúng vì trục tung biểu diễn li độ $x$, do đó giá trị tung độ chính là tọa độ li độ thực tế của vật
\item (Đúng) Độ dốc của đồ thị tại một điểm bất kì biểu thị giá trị vận tốc tức thời của vật tại thời điểm đó. (Vì): Mệnh đề thứ hai đúng vì theo định nghĩa vật lí, vận tốc là đạo hàm của li độ theo thời gian, tương ứng với độ dốc của tiếp tuyến đồ thị $x - t$
\item (Đúng) Nếu đường hình sin của đồ thị có xu hướng đi xuống cắt trục hoành thì vật đang chuyển động qua vị trí cân bằng theo chiều âm. (Vì): Mệnh đề thứ ba đúng vì khi đồ thị cắt trục hoành ($x = 0$) đi xuống nghĩa là li độ đang giảm dần về giá trị âm, tức là vật chuyển động ngược chiều dương
\item (Sai) Khoảng cách giữa hai điểm cắt trục hoành liên tiếp của đồ thị luôn có giá trị bằng một chu kì dao động $T$. (Vì): Mệnh đề thứ tư sai vì hai điểm cắt trục hoành liên tiếp tương ứng với hai lần liên tiếp vật đi qua vị trí cân bằng, khoảng thời gian này chỉ bằng một nửa chu kì $\dfrac{T}{2}$
\end{itemchoice}
}
\end{ex}

% ===== Câu 131 =====
% ID: L11C1B1-28-DS
% Mức: TH
\begin{ex}
% ID: L11C1B1-28-DS
% Mức: TH
Khi phân tích các đặc điểm của đồ thị li độ – thời gian của một vật dao động điều hòa quanh vị trí cân bằng, hãy nhận xét các mệnh đề sau:
\choiceTF
{\True Tọa độ li độ $x$ của vật tại một thời điểm được xác định trực tiếp bằng tung độ của điểm biểu diễn tương ứng trên đồ thị.}
{\True Độ dốc của đồ thị tại một điểm bất kì biểu thị giá trị vận tốc tức thời của vật tại thời điểm đó.}
{\True Nếu đường hình sin của đồ thị có xu hướng đi xuống cắt trục hoành thì vật đang chuyển động qua vị trí cân bằng theo chiều âm.}
{Khoảng cách giữa hai điểm cắt trục hoành liên tiếp của đồ thị luôn có giá trị bằng một chu kì dao động $T$.}
\loigiai{
\begin{itemchoice}
\item (Đúng) Tọa độ li độ $x$ của vật tại một thời điểm được xác định trực tiếp bằng tung độ của điểm biểu diễn tương ứng trên đồ thị. (Vì): Mệnh đề thứ nhất đúng vì trục tung biểu diễn li độ $x$, do đó giá trị tung độ chính là tọa độ li độ thực tế của vật
\item (Đúng) Độ dốc của đồ thị tại một điểm bất kì biểu thị giá trị vận tốc tức thời của vật tại thời điểm đó. (Vì): Mệnh đề thứ hai đúng vì theo định nghĩa vật lí, vận tốc là đạo hàm của li độ theo thời gian, tương ứng với độ dốc của tiếp tuyến đồ thị $x - t$
\item (Đúng) Nếu đường hình sin của đồ thị có xu hướng đi xuống cắt trục hoành thì vật đang chuyển động qua vị trí cân bằng theo chiều âm. (Vì): Mệnh đề thứ ba đúng vì khi đồ thị cắt trục hoành ($x = 0$) đi xuống nghĩa là li độ đang giảm dần về giá trị âm, tức là vật chuyển động ngược chiều dương
\item (Sai) Khoảng cách giữa hai điểm cắt trục hoành liên tiếp của đồ thị luôn có giá trị bằng một chu kì dao động $T$. (Vì): Mệnh đề thứ tư sai vì hai điểm cắt trục hoành liên tiếp tương ứng với hai lần liên tiếp vật đi qua vị trí cân bằng, khoảng thời gian này chỉ bằng một nửa chu kì $\dfrac{T}{2}$
\end{itemchoice}
}
\end{ex}

% ===== Câu 133 =====
% ID: L11C1B1-29-DS
% Mức: TH
\begin{ex}
% ID: L11C1B1-29-DS
% Mức: TH
Khi phân tích các đặc điểm của đồ thị li độ – thời gian của một vật dao động điều hòa quanh vị trí cân bằng, hãy nhận xét các mệnh đề sau:
\choiceTF
{\True Tọa độ li độ $x$ của vật tại một thời điểm được xác định trực tiếp bằng tung độ của điểm biểu diễn tương ứng trên đồ thị.}
{\True Độ dốc của đồ thị tại một điểm bất kì biểu thị giá trị vận tốc tức thời của vật tại thời điểm đó.}
{\True Nếu đường hình sin của đồ thị có xu hướng đi xuống cắt trục hoành thì vật đang chuyển động qua vị trí cân bằng theo chiều âm.}
{Khoảng cách giữa hai điểm cắt trục hoành liên tiếp của đồ thị luôn có giá trị bằng một chu kì dao động $T$.}
\loigiai{
\begin{itemchoice}
\item (Đúng) Tọa độ li độ $x$ của vật tại một thời điểm được xác định trực tiếp bằng tung độ của điểm biểu diễn tương ứng trên đồ thị. (Vì): Mệnh đề thứ nhất đúng vì trục tung biểu diễn li độ $x$, do đó giá trị tung độ chính là tọa độ li độ thực tế của vật
\item (Đúng) Độ dốc của đồ thị tại một điểm bất kì biểu thị giá trị vận tốc tức thời của vật tại thời điểm đó. (Vì): Mệnh đề thứ hai đúng vì theo định nghĩa vật lí, vận tốc là đạo hàm của li độ theo thời gian, tương ứng với độ dốc của tiếp tuyến đồ thị $x - t$
\item (Đúng) Nếu đường hình sin của đồ thị có xu hướng đi xuống cắt trục hoành thì vật đang chuyển động qua vị trí cân bằng theo chiều âm. (Vì): Mệnh đề thứ ba đúng vì khi đồ thị cắt trục hoành ($x = 0$) đi xuống nghĩa là li độ đang giảm dần về giá trị âm, tức là vật chuyển động ngược chiều dương
\item (Sai) Khoảng cách giữa hai điểm cắt trục hoành liên tiếp của đồ thị luôn có giá trị bằng một chu kì dao động $T$. (Vì): Mệnh đề thứ tư sai vì hai điểm cắt trục hoành liên tiếp tương ứng với hai lần liên tiếp vật đi qua vị trí cân bằng, khoảng thời gian này chỉ bằng một nửa chu kì $\dfrac{T}{2}$
\end{itemchoice}
}
\end{ex}

% ===== Câu 135 =====
% ID: L11C1B1-31-DS
% Mức: VD
\begin{ex}
% ID: L11C1B1-31-DS
% Mức: VD
Một vật dao động điều hòa có đồ thị li độ – thời gian dạng hình sin. Tại thời điểm $t = 0$, đồ thị bắt đầu từ biên dương $x_0 = 5$ cm và đi xuống vị trí cân bằng lần đầu tiên tại thời điểm $t = 0,1$ s. Xét các thông số của dao động:
\choiceTF
{\True Biên độ dao động của vật bằng $5$ cm.}
{\True Chu kì dao động của vật bằng $0,4$ s.}
{\True Tần số góc của dao động đạt giá trị $5\pi$ rad/s.}
{Phương trình dao động điều hòa của vật là $x = 5\cos\left(5\pi t + \dfrac{\pi}{2}\right)$ (cm).}
\loigiai{
\begin{itemchoice}
\item (Đúng) Biên độ dao động của vật bằng $5$ cm. (Vì): (Đúng) Biên độ dao động của vật bằng $5$ cm
\item (Đúng) Chu kì dao động của vật bằng $0,4$ s. (Vì): (Đúng) Chu kì dao động của vật bằng $0,4$ s
\item (Đúng) Tần số góc của dao động đạt giá trị $5\pi$ rad/s. (Vì): m
\item (Sai) Phương trình dao động điều hòa của vật là $x = 5\cos\left(5\pi t + \dfrac{\pi}{2}\right)$ (cm). (Vì): (Sai) Phương trình dao động điều hòa của vật là $x = 5\cos\left(5\pi t + \dfrac{\pi}{2}\right)$
\end{itemchoice}
}
\end{ex}

% ===== Câu 137 =====
% ID: L11C1B1-32-DS
% Mức: VD
\begin{ex}
% ID: L11C1B1-32-DS
% Mức: VD
Một vật dao động điều hòa có đồ thị li độ – thời gian dạng hình sin. Tại thời điểm $t = 0$, đồ thị bắt đầu từ biên dương $x_0 = 5$ cm và đi xuống vị trí cân bằng lần đầu tiên tại thời điểm $t = 0,1$ s. Xét các thông số của dao động:
\choiceTF
{\True Biên độ dao động của vật bằng $5$ cm.}
{\True Chu kì dao động của vật bằng $0,4$ s.}
{\True Tần số góc của dao động đạt giá trị $5\pi$ rad/s.}
{Phương trình dao động điều hòa của vật là $x = 5\cos\left(5\pi t + \dfrac{\pi}{2}\right)$ (cm).}
\loigiai{
\begin{itemchoice}
\item (Đúng) Biên độ dao động của vật bằng $5$ cm. (Vì): (Đúng) Biên độ dao động của vật bằng $5$ cm
\item (Đúng) Chu kì dao động của vật bằng $0,4$ s. (Vì): (Đúng) Chu kì dao động của vật bằng $0,4$ s
\item (Đúng) Tần số góc của dao động đạt giá trị $5\pi$ rad/s. (Vì): m
\item (Sai) Phương trình dao động điều hòa của vật là $x = 5\cos\left(5\pi t + \dfrac{\pi}{2}\right)$ (cm). (Vì): (Sai) Phương trình dao động điều hòa của vật là $x = 5\cos\left(5\pi t + \dfrac{\pi}{2}\right)$
\end{itemchoice}
}
\end{ex}

% ===== Câu 139 =====
% ID: L11C1B1-34-DS
% Mức: VD
\begin{ex}
% ID: L11C1B1-34-DS
% Mức: VD
Một vật dao động điều hòa có đồ thị li độ – thời gian dạng hình sin. Tại thời điểm $t = 0$, đồ thị bắt đầu từ biên dương $x_0 = 5$ cm và đi xuống vị trí cân bằng lần đầu tiên tại thời điểm $t = 0,1$ s. Xét các thông số của dao động:
\choiceTF
{\True Biên độ dao động của vật bằng $5$ cm.}
{\True Chu kì dao động của vật bằng $0,4$ s.}
{\True Tần số góc của dao động đạt giá trị $5\pi$ rad/s.}
{Phương trình dao động điều hòa của vật là $x = 5\cos\left(5\pi t + \dfrac{\pi}{2}\right)$ (cm).}
\loigiai{
\begin{itemchoice}
\item (Đúng) Biên độ dao động của vật bằng $5$ cm. (Vì): (Đúng) Biên độ dao động của vật bằng $5$ cm
\item (Đúng) Chu kì dao động của vật bằng $0,4$ s. (Vì): (Đúng) Chu kì dao động của vật bằng $0,4$ s
\item (Đúng) Tần số góc của dao động đạt giá trị $5\pi$ rad/s. (Vì): m
\item (Sai) Phương trình dao động điều hòa của vật là $x = 5\cos\left(5\pi t + \dfrac{\pi}{2}\right)$ (cm). (Vì): (Sai) Phương trình dao động điều hòa của vật là $x = 5\cos\left(5\pi t + \dfrac{\pi}{2}\right)$
\end{itemchoice}
}
\end{ex}

% ===== Câu 141 =====
% ID: L11C1B1-35-DS
% Mức: VD
\begin{ex}
% ID: L11C1B1-35-DS
% Mức: VD
Quan sát đồ thị dao động điều hòa của một vật có đường hình sin biểu diễn trên hệ tọa độ $x - t$. Tại thời điểm $t = 0$, đồ thị đi qua gốc tọa độ và đang đi lên theo chiều dương, sau đó đạt tới biên dương với tọa độ li độ $10$ cm lần đầu tiên tại mốc thời gian $t = 0,25$ s. Khảo sát các mệnh đề sau:
\choiceTF
{\True Biên độ dao động của vật là $A = 10$ cm.}
{\True Chu kì dao động của vật là $T = 1,0$ s.}
{Pha ban đầu của dao động bằng $\varphi = \dfrac{\pi}{2}$ rad.}
{\True Phương trình li độ của vật là $x = 10\cos\left(2\pi t - \dfrac{\pi}{2}\right)$ (cm).}
\loigiai{
\begin{itemchoice}
\item (Đúng) Biên độ dao động của vật là $A = 10$ cm. (Vì): Mệnh đề thứ nhất đúng vì tọa độ đỉnh của đường hình sin trên đồ thị tương ứng với biên độ $A = 10$ cm
\item (Đúng) Chu kì dao động của vật là $T = 1,0$ s. (Vì): Mệnh đề thứ hai đúng vì khoảng thời gian vật đi từ vị trí cân bằng đến vị trí biên dương lần đầu tiên bằng một phần tư chu kì: $\dfrac{T}{4} = 0,25 \text{ s} \Rightarrow T = 1,0$ s
\item (Sai) Pha ban đầu của dao động bằng $\varphi = \dfrac{\pi}{2}$ rad. (Vì): Mệnh đề thứ ba sai vì tại gốc thời gian vật đi qua vị trí cân bằng theo chiều dương nên pha ban đầu có giá trị âm: $\varphi = -\dfrac{\pi}{2}$ rad
\item (Đúng) Phương trình li độ của vật là $x = 10\cos\left(2\pi t - \dfrac{\pi}{2}\right)$ (cm). (Vì): (Đúng) Phương trình li độ của vật là $x = 10\cos\left(2\pi t - \dfrac{\pi}{2}\right)$
\end{itemchoice}
}
\end{ex}

% ===== Câu 143 =====
% ID: L11C1B1-36-DS
% Mức: VD
\begin{ex}
% ID: L11C1B1-36-DS
% Mức: VD
Quan sát đồ thị dao động điều hòa của một vật có đường hình sin biểu diễn trên hệ tọa độ $x - t$. Tại thời điểm $t = 0$, đồ thị đi qua gốc tọa độ và đang đi lên theo chiều dương, sau đó đạt tới biên dương với tọa độ li độ $10$ cm lần đầu tiên tại mốc thời gian $t = 0,25$ s. Khảo sát các mệnh đề sau:
\choiceTF
{\True Biên độ dao động của vật là $A = 10$ cm.}
{\True Chu kì dao động của vật là $T = 1,0$ s.}
{Pha ban đầu của dao động bằng $\varphi = \dfrac{\pi}{2}$ rad.}
{\True Phương trình li độ của vật là $x = 10\cos\left(2\pi t - \dfrac{\pi}{2}\right)$ (cm).}
\loigiai{
\begin{itemchoice}
\item (Đúng) Biên độ dao động của vật là $A = 10$ cm. (Vì): Mệnh đề thứ nhất đúng vì tọa độ đỉnh của đường hình sin trên đồ thị tương ứng với biên độ $A = 10$ cm
\item (Đúng) Chu kì dao động của vật là $T = 1,0$ s. (Vì): Mệnh đề thứ hai đúng vì khoảng thời gian vật đi từ vị trí cân bằng đến vị trí biên dương lần đầu tiên bằng một phần tư chu kì: $\dfrac{T}{4} = 0,25 \text{ s} \Rightarrow T = 1,0$ s
\item (Sai) Pha ban đầu của dao động bằng $\varphi = \dfrac{\pi}{2}$ rad. (Vì): Mệnh đề thứ ba sai vì tại gốc thời gian vật đi qua vị trí cân bằng theo chiều dương nên pha ban đầu có giá trị âm: $\varphi = -\dfrac{\pi}{2}$ rad
\item (Đúng) Phương trình li độ của vật là $x = 10\cos\left(2\pi t - \dfrac{\pi}{2}\right)$ (cm). (Vì): (Đúng) Phương trình li độ của vật là $x = 10\cos\left(2\pi t - \dfrac{\pi}{2}\right)$
\end{itemchoice}
}
\end{ex}

% ===== Câu 145 =====
% ID: L11C1B1-38-DS
% Mức: VD
\begin{ex}
% ID: L11C1B1-38-DS
% Mức: VD
Quan sát đồ thị dao động điều hòa của một vật có đường hình sin biểu diễn trên hệ tọa độ $x - t$. Tại thời điểm $t = 0$, đồ thị đi qua gốc tọa độ và đang đi lên theo chiều dương, sau đó đạt tới biên dương với tọa độ li độ $10$ cm lần đầu tiên tại mốc thời gian $t = 0,25$ s. Khảo sát các mệnh đề sau:
\choiceTF
{\True Biên độ dao động của vật là $A = 10$ cm.}
{\True Chu kì dao động của vật là $T = 1,0$ s.}
{Pha ban đầu của dao động bằng $\varphi = \dfrac{\pi}{2}$ rad.}
{\True Phương trình li độ của vật là $x = 10\cos\left(2\pi t - \dfrac{\pi}{2}\right)$ (cm).}
\loigiai{
\begin{itemchoice}
\item (Đúng) Biên độ dao động của vật là $A = 10$ cm. (Vì): Mệnh đề thứ nhất đúng vì tọa độ đỉnh của đường hình sin trên đồ thị tương ứng với biên độ $A = 10$ cm
\item (Đúng) Chu kì dao động của vật là $T = 1,0$ s. (Vì): Mệnh đề thứ hai đúng vì khoảng thời gian vật đi từ vị trí cân bằng đến vị trí biên dương lần đầu tiên bằng một phần tư chu kì: $\dfrac{T}{4} = 0,25 \text{ s} \Rightarrow T = 1,0$ s
\item (Sai) Pha ban đầu của dao động bằng $\varphi = \dfrac{\pi}{2}$ rad. (Vì): Mệnh đề thứ ba sai vì tại gốc thời gian vật đi qua vị trí cân bằng theo chiều dương nên pha ban đầu có giá trị âm: $\varphi = -\dfrac{\pi}{2}$ rad
\item (Đúng) Phương trình li độ của vật là $x = 10\cos\left(2\pi t - \dfrac{\pi}{2}\right)$ (cm). (Vì): (Đúng) Phương trình li độ của vật là $x = 10\cos\left(2\pi t - \dfrac{\pi}{2}\right)$
\end{itemchoice}
}
\end{ex}

% ===== Câu 147 =====
% ID: L11C1B1-40-DS
% Mức: VD
\begin{ex}
% ID: L11C1B1-40-DS
% Mức: VD
Một vật dao động điều hòa có đồ thị li độ – thời gian dạng hình sin biến thiên tuần hoàn. Khoảng thời gian đo được giữa hai đỉnh cực đại liên tiếp trên đồ thị là $0,8$ s, và khoảng cách thẳng đứng từ đỉnh cao nhất đến đáy thấp nhất của đồ thị là $16$ cm. Đánh giá các mệnh đề sau:
\choiceTF
{Biên độ dao động của vật là $16$ cm.}
{\True Chu kì dao động của vật bằng $0,8$ s.}
{\True Tần số dao động của vật bằng $1,25$ Hz.}
{\True Tần số góc của dao động có giá trị bằng $2,5\pi$ rad/s.}
\loigiai{
\begin{itemchoice}
\item (Sai) Biên độ dao động của vật là $16$ cm. (Vì): Mệnh đề thứ nhất sai vì khoảng cách thẳng đứng từ đỉnh đến đáy đồ thị là $2A = 16$ cm, suy ra biên độ dao động của vật là $A = 8$ cm
\item (Đúng) Chu kì dao động của vật bằng $0,8$ s. (Vì): Mệnh đề thứ hai đúng vì khoảng thời gian giữa hai đỉnh cực đại (hai lần liên tiếp vật đạt biên dương) chính là một chu kì dao động $T = 0,8$ s
\item (Đúng) Tần số dao động của vật bằng $1,25$ Hz. (Vì): Mệnh đề thứ ba đúng vì tần số dao động là: $f = \dfrac{1}{T} = \dfrac{1}{0,8} = 1,25$ Hz
\item (Đúng) Tần số góc của dao động có giá trị bằng $2,5\pi$ rad/s. (Vì): Mệnh đề thứ tư đúng vì tần số góc của dao động là: $\omega = 2\pi f = 2\pi \cdot 1,25 = 2,5\pi$ rad/s
\end{itemchoice}
}
\end{ex}

% ===== Câu 149 =====
% ID: L11C1B1-41-DS
% Mức: VD
\begin{ex}
% ID: L11C1B1-41-DS
% Mức: VD
Một vật dao động điều hòa có đồ thị li độ – thời gian dạng hình sin biến thiên tuần hoàn. Khoảng thời gian đo được giữa hai đỉnh cực đại liên tiếp trên đồ thị là $0,8$ s, và khoảng cách thẳng đứng từ đỉnh cao nhất đến đáy thấp nhất của đồ thị là $16$ cm. Đánh giá các mệnh đề sau:
\choiceTF
{Biên độ dao động của vật là $16$ cm.}
{\True Chu kì dao động của vật bằng $0,8$ s.}
{\True Tần số dao động của vật bằng $1,25$ Hz.}
{\True Tần số góc của dao động có giá trị bằng $2,5\pi$ rad/s.}
\loigiai{
\begin{itemchoice}
\item (Sai) Biên độ dao động của vật là $16$ cm. (Vì): Mệnh đề thứ nhất sai vì khoảng cách thẳng đứng từ đỉnh đến đáy đồ thị là $2A = 16$ cm, suy ra biên độ dao động của vật là $A = 8$ cm
\item (Đúng) Chu kì dao động của vật bằng $0,8$ s. (Vì): Mệnh đề thứ hai đúng vì khoảng thời gian giữa hai đỉnh cực đại (hai lần liên tiếp vật đạt biên dương) chính là một chu kì dao động $T = 0,8$ s
\item (Đúng) Tần số dao động của vật bằng $1,25$ Hz. (Vì): Mệnh đề thứ ba đúng vì tần số dao động là: $f = \dfrac{1}{T} = \dfrac{1}{0,8} = 1,25$ Hz
\item (Đúng) Tần số góc của dao động có giá trị bằng $2,5\pi$ rad/s. (Vì): Mệnh đề thứ tư đúng vì tần số góc của dao động là: $\omega = 2\pi f = 2\pi \cdot 1,25 = 2,5\pi$ rad/s
\end{itemchoice}
}
\end{ex}

% ===== Câu 151 =====
% ID: L11C1B1-42-DS
% Mức: VD
\begin{ex}
% ID: L11C1B1-42-DS
% Mức: VD
Một vật dao động điều hòa có đồ thị li độ – thời gian dạng hình sin biến thiên tuần hoàn. Khoảng thời gian đo được giữa hai đỉnh cực đại liên tiếp trên đồ thị là $0,8$ s, và khoảng cách thẳng đứng từ đỉnh cao nhất đến đáy thấp nhất của đồ thị là $16$ cm. Đánh giá các mệnh đề sau:
\choiceTF
{Biên độ dao động của vật là $16$ cm.}
{\True Chu kì dao động của vật bằng $0,8$ s.}
{\True Tần số dao động của vật bằng $1,25$ Hz.}
{\True Tần số góc của dao động có giá trị bằng $2,5\pi$ rad/s.}
\loigiai{
\begin{itemchoice}
\item (Sai) Biên độ dao động của vật là $16$ cm. (Vì): Mệnh đề thứ nhất sai vì khoảng cách thẳng đứng từ đỉnh đến đáy đồ thị là $2A = 16$ cm, suy ra biên độ dao động của vật là $A = 8$ cm
\item (Đúng) Chu kì dao động của vật bằng $0,8$ s. (Vì): Mệnh đề thứ hai đúng vì khoảng thời gian giữa hai đỉnh cực đại (hai lần liên tiếp vật đạt biên dương) chính là một chu kì dao động $T = 0,8$ s
\item (Đúng) Tần số dao động của vật bằng $1,25$ Hz. (Vì): Mệnh đề thứ ba đúng vì tần số dao động là: $f = \dfrac{1}{T} = \dfrac{1}{0,8} = 1,25$ Hz
\item (Đúng) Tần số góc của dao động có giá trị bằng $2,5\pi$ rad/s. (Vì): Mệnh đề thứ tư đúng vì tần số góc của dao động là: $\omega = 2\pi f = 2\pi \cdot 1,25 = 2,5\pi$ rad/s
\end{itemchoice}
}
\end{ex}

% ===== Câu 153 =====
% ID: L11C1B1-43-DS
% Mức: VD
\begin{ex}
% ID: L11C1B1-43-DS
% Mức: VD
Một vật dao động điều hòa có đồ thị li độ – thời gian hình sin. Tại thời điểm $t = 0$, vật ở vị trí biên âm $x_0 = -4$ cm. Đường biểu diễn đi qua vị trí cân bằng lần thứ hai tại thời điểm mốc $t = 0,75$ s. Khảo sát các mệnh đề sau:
\choiceTF
{\True Biên độ dao động của vật bằng $4$ cm.}
{\True Kể từ thời điểm ban đầu, vật đi qua vị trí cân bằng lần thứ nhất tại thời điểm $t = 0,25$ s.}
{\True Chu kì dao động của vật là $1,0$ s.}
{Pha ban đầu của dao động bằng $0$ rad.}
\loigiai{
\begin{itemchoice}
\item (Đúng) Biên độ dao động của vật bằng $4$ cm. (Vì): Mệnh đề thứ nhất đúng vì vật xuất phát từ biên âm nên khoảng cách từ vị trí cân bằng tới biên là $A = 4$ cm
\item (Đúng) Kể từ thời điểm ban đầu, vật đi qua vị trí cân bằng lần thứ nhất tại thời điểm $t = 0,25$ s. (Vì): Mệnh đề thứ hai đúng vì thời gian vật đi từ biên âm đến vị trí cân bằng lần đầu tiên là một phần tư chu kì: $t_1 = \dfrac{T}{4}$
\item (Đúng) Chu kì dao động của vật là $1,0$ s. (Vì): Mệnh đề thứ ba đúng vì thời gian từ lúc xuất phát ở biên âm đến lúc qua vị trí cân bằng lần thứ hai là ba phần tư chu kì: $\dfrac{3T}{4} = 0,75 \text{ s} \Rightarrow T = 1,0$ s. Do đó $t_1 = \dfrac{1,0}{4} = 0,25$ s
\item (Sai) Pha ban đầu của dao động bằng $0$ rad. (Vì): Mệnh đề thứ tư sai vì ban đầu vật ở vị trí biên âm nên pha ban đầu phải bằng $\varphi = \pi$ rad
\end{itemchoice}
}
\end{ex}

% ===== Câu 155 =====
% ID: L11C1B1-44-DS
% Mức: VD
\begin{ex}
% ID: L11C1B1-44-DS
% Mức: VD
Một vật dao động điều hòa có đồ thị li độ – thời gian hình sin. Tại thời điểm $t = 0$, vật ở vị trí biên âm $x_0 = -4$ cm. Đường biểu diễn đi qua vị trí cân bằng lần thứ hai tại thời điểm mốc $t = 0,75$ s. Khảo sát các mệnh đề sau:
\choiceTF
{\True Biên độ dao động của vật bằng $4$ cm.}
{\True Kể từ thời điểm ban đầu, vật đi qua vị trí cân bằng lần thứ nhất tại thời điểm $t = 0,25$ s.}
{\True Chu kì dao động của vật là $1,0$ s.}
{Pha ban đầu của dao động bằng $0$ rad.}
\loigiai{
\begin{itemchoice}
\item (Đúng) Biên độ dao động của vật bằng $4$ cm. (Vì): Mệnh đề thứ nhất đúng vì vật xuất phát từ biên âm nên khoảng cách từ vị trí cân bằng tới biên là $A = 4$ cm
\item (Đúng) Kể từ thời điểm ban đầu, vật đi qua vị trí cân bằng lần thứ nhất tại thời điểm $t = 0,25$ s. (Vì): Mệnh đề thứ hai đúng vì thời gian vật đi từ biên âm đến vị trí cân bằng lần đầu tiên là một phần tư chu kì: $t_1 = \dfrac{T}{4}$
\item (Đúng) Chu kì dao động của vật là $1,0$ s. (Vì): Mệnh đề thứ ba đúng vì thời gian từ lúc xuất phát ở biên âm đến lúc qua vị trí cân bằng lần thứ hai là ba phần tư chu kì: $\dfrac{3T}{4} = 0,75 \text{ s} \Rightarrow T = 1,0$ s. Do đó $t_1 = \dfrac{1,0}{4} = 0,25$ s
\item (Sai) Pha ban đầu của dao động bằng $0$ rad. (Vì): Mệnh đề thứ tư sai vì ban đầu vật ở vị trí biên âm nên pha ban đầu phải bằng $\varphi = \pi$ rad
\end{itemchoice}
}
\end{ex}

% ===== Câu 157 =====
% ID: L11C1B1-45-DS
% Mức: VD
\begin{ex}
% ID: L11C1B1-45-DS
% Mức: VD
Một vật dao động điều hòa có đồ thị li độ – thời gian hình sin. Tại thời điểm $t = 0$, vật ở vị trí biên âm $x_0 = -4$ cm. Đường biểu diễn đi qua vị trí cân bằng lần thứ hai tại thời điểm mốc $t = 0,75$ s. Khảo sát các mệnh đề sau:
\choiceTF
{\True Biên độ dao động của vật bằng $4$ cm.}
{\True Kể từ thời điểm ban đầu, vật đi qua vị trí cân bằng lần thứ nhất tại thời điểm $t = 0,25$ s.}
{\True Chu kì dao động của vật là $1,0$ s.}
{Pha ban đầu của dao động bằng $0$ rad.}
\loigiai{
\begin{itemchoice}
\item (Đúng) Biên độ dao động của vật bằng $4$ cm. (Vì): Mệnh đề thứ nhất đúng vì vật xuất phát từ biên âm nên khoảng cách từ vị trí cân bằng tới biên là $A = 4$ cm
\item (Đúng) Kể từ thời điểm ban đầu, vật đi qua vị trí cân bằng lần thứ nhất tại thời điểm $t = 0,25$ s. (Vì): Mệnh đề thứ hai đúng vì thời gian vật đi từ biên âm đến vị trí cân bằng lần đầu tiên là một phần tư chu kì: $t_1 = \dfrac{T}{4}$
\item (Đúng) Chu kì dao động của vật là $1,0$ s. (Vì): Mệnh đề thứ ba đúng vì thời gian từ lúc xuất phát ở biên âm đến lúc qua vị trí cân bằng lần thứ hai là ba phần tư chu kì: $\dfrac{3T}{4} = 0,75 \text{ s} \Rightarrow T = 1,0$ s. Do đó $t_1 = \dfrac{1,0}{4} = 0,25$ s
\item (Sai) Pha ban đầu của dao động bằng $0$ rad. (Vì): Mệnh đề thứ tư sai vì ban đầu vật ở vị trí biên âm nên pha ban đầu phải bằng $\varphi = \pi$ rad
\end{itemchoice}
}
\end{ex}

% ===== Câu 170 =====
% ID: L11C1B1-46-TN
% Mức: TH
\begin{ex}
% ID: L11C1B1-46-TN
% Mức: TH
Từ đồ thị li độ - thời gian, ta thấy tại $t = 0$, li độ là $x = -3$ cm; vật đạt biên độ âm $-6$ cm tại $t = 0{,}1$ s; quay trở về li độ $-3$ cm tại $t = 0{,}2$ s. Biên độ dao động là bao nhiêu?
\choice
{$A = 3$ cm}
{$A = 6$ cm}
{\True $A = 6$ cm}
{$A = 9$ cm}
\loigiai{
Biên độ là độ lệch cực đại từ vị trí cân bằng. Từ đồ thị, vật đạt giá trị cực đại là $\pm 6$ cm, do đó biên độ $A = 6$ cm.
}
\end{ex}

% ===== Câu 172 =====
% ID: L11C1B1-47-TN
% Mức: TH
\begin{ex}
% ID: L11C1B1-47-TN
% Mức: TH
Từ đồ thị li độ - thời gian, ta thấy tại $t = 0$, li độ là $x = -3$ cm; vật đạt biên độ âm $-6$ cm tại $t = 0{,}1$ s; quay trở về li độ $-3$ cm tại $t = 0{,}2$ s. Biên độ dao động là bao nhiêu?
\choice
{$A = 3$ cm}
{$A = 6$ cm}
{\True $A = 6$ cm}
{$A = 9$ cm}
\loigiai{
Biên độ là độ lệch cực đại từ vị trí cân bằng. Từ đồ thị, vật đạt giá trị cực đại là $\pm 6$ cm, do đó biên độ $A = 6$ cm.
}
\end{ex}

% ===== Câu 174 =====
% ID: L11C1B1-48-TN
% Mức: TH
\begin{ex}
% ID: L11C1B1-48-TN
% Mức: TH
Từ đồ thị li độ - thời gian, ta thấy tại $t = 0$, li độ là $x = -3$ cm; vật đạt biên độ âm $-6$ cm tại $t = 0{,}1$ s; quay trở về li độ $-3$ cm tại $t = 0{,}2$ s. Biên độ dao động là bao nhiêu?
\choice
{$A = 3$ cm}
{$A = 6$ cm}
{\True $A = 6$ cm}
{$A = 9$ cm}
\loigiai{
Biên độ là độ lệch cực đại từ vị trí cân bằng. Từ đồ thị, vật đạt giá trị cực đại là $\pm 6$ cm, do đó biên độ $A = 6$ cm.
}
\end{ex}

% ===== Câu 176 =====
% ID: L11C1B1-49-TN
% Mức: TH
\begin{ex}
% ID: L11C1B1-49-TN
% Mức: TH
Đồ thị li độ - thời gian cho thấy biên độ dao động là $6\,\text{cm}$ và trong khoảng thời gian từ $t = 0$ đến $t = 0,5$ s, vật hoàn thành một dao động toàn phần. Tần số dao động là bao nhiêu?
\choice
{$f = 0{,}5$ Hz}
{$f = 1$ Hz}
{\True $f = 2$ Hz}
{$f = 4$ Hz}
\loigiai{
Từ đồ thị li độ - thời gian, ta có khoảng thời gian vật hoàn thành một dao động toàn phần chính là chu kì dao động:
$T=0,5\,\text{s}$
Tần số dao động của vật là:
$f = \frac{1}{T} = \frac{1}{0,5} = 2\,\text{Hz}$
Do đó, ta chọn đáp án C.
}
\end{ex}

% ===== Câu 178 =====
% ID: L11C1B1-50-TN
% Mức: TH
\begin{ex}
% ID: L11C1B1-50-TN
% Mức: TH
Đồ thị li độ - thời gian cho thấy biên độ dao động là $6\,\text{cm}$ và trong khoảng thời gian từ $t = 0$ đến $t = 0,5$ s, vật hoàn thành một dao động toàn phần. Tần số dao động là bao nhiêu?
\choice
{$f = 0{,}5$ Hz}
{$f = 1$ Hz}
{\True $f = 2$ Hz}
{$f = 4$ Hz}
\loigiai{
Tần số là số dao động hoàn thành trong một đơn vị thời gian.
 Nếu một dao động toàn phần mất $0,5\,\text{s}$, tức là $T = 0,5$ s, thì:
  \begin{aligned}
 f &= $\frac{1}{T}$ &= $\frac{1}{0,5}$ &= 2 $\\text{Hz}$.
 \end{aligned} 
 Vậy tần số dao động là $2\,\text{Hz}$.
}
\end{ex}

% ===== Câu 180 =====
% ID: L11C1B1-51-TN
% Mức: TH
\begin{ex}
% ID: L11C1B1-51-TN
% Mức: TH
Đồ thị li độ - thời gian cho thấy biên độ dao động là $6\,\text{cm}$ và trong khoảng thời gian từ $t = 0$ đến $t = 0,5$ s, vật hoàn thành một dao động toàn phần. Tần số dao động là bao nhiêu?
\choice
{$f = 0{,}5$ Hz}
{$f = 1$ Hz}
{\True $f = 2$ Hz}
{$f = 4$ Hz}
\loigiai{
Tần số là số dao động hoàn thành trong một đơn vị thời gian.
 Nếu một dao động toàn phần mất $0,5\,\text{s}$, tức là $T = 0,5$ s, thì:
  \begin{aligned}
 f &= $\frac{1}{T}$ &= $\frac{1}{0,5}$ &= 2 $\\text{Hz}$.
 \end{aligned} 
 Vậy tần số dao động là $2\,\text{Hz}$.
}
\end{ex}

% ===== Câu 182 =====
% ID: L11C1B1-52-TN
% Mức: TH
\begin{ex}
% ID: L11C1B1-52-TN
% Mức: TH
Từ đồ thị li độ - thời gian của một dao động điều hòa, ta thấy tại $t = 0$, li độ là $0\,\text{cm}$; vật đạt biên độ dương $4\,\text{cm}$ tại $t = 0,1$ s; quay trở về li độ 0 tại $t = 0,2$ s. Chu kì của dao động là bao nhiêu?
\choice
{$T = 0{,}1$ s}
{$T = 0{,}2$ s}
{\True $T = 0{,}4$ s}
{$T = 0{,}5$ s}
\loigiai{
Từ đồ thị, vật đi từ li độ 0 đến biên dương $4\,\text{cm}$ mất $0,1\,\text{s}$, sau đó quay về li độ 0 tại t = $0,2\,\text{s}$.
Điều này cho thấy $0,2\,\text{s}$ là thời gian để vật hoàn thành nửa chu kì dao động.
Ta có:
$\frac{T}{2} = 0,2\,\text{s}\Rightarrow T = 2 \cdot 0,2 = 0,4\,\text{s}$
}
\end{ex}

% ===== Câu 184 =====
% ID: L11C1B1-53-TN
% Mức: TH
\begin{ex}
% ID: L11C1B1-53-TN
% Mức: TH
Từ đồ thị li độ - thời gian của một dao động điều hòa, ta thấy tại $t = 0$, li độ là $0\,\text{cm}$; vật đạt biên độ dương $4\,\text{cm}$ tại $t = 0,1$ s; quay trở về li độ 0 tại $t = 0,2$ s. Chu kì của dao động là bao nhiêu?
\choice
{$T = 0{,}1$ s}
{$T = 0{,}2$ s}
{\True $T = 0{,}4$ s}
{$T = 0{,}5$ s}
\loigiai{
Từ đồ thị, vật đi từ li độ 0 đến biên dương $4\,\text{cm}$ mất $0,1\,\text{s}$, sau đó quay về li độ 0 trong $0,1\,\text{s}$ nữa, tổng cộng $0,2\,\text{s}$ để hoàn thành nửa chu kì.
 Do đó:
  \begin{aligned}
 $\frac{T}{2}$ &= 0,2 $\\text{s}T$ &= 0,4 $\\text{s}$.
 \end{aligned} 
 Vậy chu kì dao động là $0,4\,\text{s}$.
}
\end{ex}

% ===== Câu 186 =====
% ID: L11C1B1-54-TN
% Mức: TH
\begin{ex}
% ID: L11C1B1-54-TN
% Mức: TH
Từ đồ thị li độ - thời gian của một dao động điều hòa, ta thấy tại $t = 0$, li độ là $0\,\text{cm}$; vật đạt biên độ dương $4\,\text{cm}$ tại $t = 0,1$ s; quay trở về li độ 0 tại $t = 0,2$ s. Chu kì của dao động là bao nhiêu?
\choice
{$T = 0{,}1$ s}
{$T = 0{,}2$ s}
{\True $T = 0{,}4$ s}
{$T = 0{,}5$ s}
\loigiai{
Từ đồ thị, vật đi từ li độ 0 đến biên dương $4\,\text{cm}$ mất $0,1\,\text{s}$, sau đó quay về li độ 0 trong $0,1\,\text{s}$ nữa, tổng cộng $0,2\,\text{s}$ để hoàn thành nửa chu kì.
 Do đó:
  \begin{aligned}
 $\frac{T}{2}$ &= 0,2 $\\text{s}T$ &= 0,4 $\\text{s}$.
 \end{aligned} 
 Vậy chu kì dao động là $0,4\,\text{s}$.
}
\end{ex}

% ===== Câu 188 =====
% ID: L11C1B1-55-TN
% Mức: VD
\begin{ex}
% ID: L11C1B1-55-TN
% Mức: VD
Từ đồ thị li độ - thời gian, biên độ dao động là $8\,\text{cm}$, lần lượt đạt các vị trí: li độ $8\,\text{cm}$ tại $t = 0$; li độ $-8$ cm tại $t = 0,2$ s; li độ $8\,\text{cm}$ tại $t = 0,4$ s. Viết phương trình li độ của dao động.
\choice
{$x = 8\cos(5\pi t)$ cm}
{$x = 8\sin(5\pi t)$ cm}
{\True $x = 8\cos(5\pi t)$ cm}
{$x = 8\cos(10\pi t)$ cm}
\loigiai{
Từ đồ thị, biên độ dao động là $A = 8\,\text{cm}$. Chu kì dao động là $T = 0,4\,\text{s}$ vì tại $t=0$ vật ở vị trí biên dương và tại $t=0,4\,\text{s}$ vật lại trở về vị trí biên dương, hoàn thành một chu kì.
Tần số góc là $\omega = \frac{2\pi}{T} = \frac{2\pi}{0,4} = 5\pi\,\text{rad/s}$.
Tại thời điểm $t = 0$, li độ $x = 8\,\text{cm}$, đây là vị trí biên dương nên pha ban đầu $\varphi_0 = 0$.
Vậy phương trình li độ của dao động là $x = A\cos(\omega t + \varphi_0) = 8\cos(5\pi t + 0) = 8\cos(5\pi t)\,\text{cm}$.
}
\end{ex}

% ===== Câu 190 =====
% ID: L11C1B1-56-TN
% Mức: VD
\begin{ex}
% ID: L11C1B1-56-TN
% Mức: VD
Từ đồ thị li độ - thời gian, biên độ dao động là $8\,\text{cm}$, lần lượt đạt các vị trí: li độ $8\,\text{cm}$ tại $t = 0$; li độ $-8$ cm tại $t = 0,2$ s; li độ $8\,\text{cm}$ tại $t = 0,4$ s. Viết phương trình li độ của dao động.
\choice
{$x = 8\cos(5\pi t)$ cm}
{$x = 8\sin(5\pi t)$ cm}
{\True $x = 8\cos(5\pi t)$ cm}
{$x = 8\cos(10\pi t)$ cm}
\loigiai{
Từ đồ thị: biên độ $A = 8$ cm; chu kì là $T = 0,4$ s (vì từ $t = 0$ đến $t = 0,4$ s hoàn thành một dao động toàn phần).
 Tần số góc:
  \begin{aligned}
 $\omega$ &= $\frac{2\pi}{T}$ &= $\frac{2\pi}{0,4}$ &= 5 $\pi\\text{rad/s}$.
 \end{aligned} 
 Tại $t = 0$, li độ $x = 8$ cm (biên độ dương), nên pha ban đầu $\varphi_0 = 0$.
 Phương trình li độ là $x = 8\cos(5\pi t)$ cm.
}
\end{ex}

% ===== Câu 192 =====
% ID: L11C1B1-57-TN
% Mức: VD
\begin{ex}
% ID: L11C1B1-57-TN
% Mức: VD
Từ đồ thị li độ - thời gian, biên độ dao động là $8\,\text{cm}$, lần lượt đạt các vị trí: li độ $8\,\text{cm}$ tại $t = 0$; li độ $-8$ cm tại $t = 0,2$ s; li độ $8\,\text{cm}$ tại $t = 0,4$ s. Viết phương trình li độ của dao động.
\choice
{$x = 8\cos(5\pi t)$ cm}
{$x = 8\sin(5\pi t)$ cm}
{\True $x = 8\cos(5\pi t)$ cm}
{$x = 8\cos(10\pi t)$ cm}
\loigiai{
Từ đồ thị: biên độ $A = 8$ cm; chu kì là $T = 0,4$ s (vì từ $t = 0$ đến $t = 0,4$ s hoàn thành một dao động toàn phần).
 Tần số góc:
  \begin{aligned}
 $\omega$ &= $\frac{2\pi}{T}$ &= $\frac{2\pi}{0,4}$ &= 5 $\pi\\text{rad/s}$.
 \end{aligned} 
 Tại $t = 0$, li độ $x = 8$ cm (biên độ dương), nên pha ban đầu $\varphi_0 = 0$.
 Phương trình li độ là $x = 8\cos(5\pi t)$ cm.
}
\end{ex}

% ===== Câu 194 =====
% ID: L11C1B1-58-TN
% Mức: VD
\begin{ex}
% ID: L11C1B1-58-TN
% Mức: VD
Đồ thị li độ - thời gian của một dao động điều hòa cho thấy vật đạt li độ cực đại $5\,\text{cm}$ tại $t = 0$, đạt li độ $0$ cm lần đầu tiên tại $t = 0,05$ s. Phương trình li độ là bao nhiêu?
\choice
{$x = 5\sin(10\pi t)$ cm}
{$x = 5\cos(10\pi t + \frac{\pi}{2})$ cm}
{\True $x = 5\cos(10\pi t)$ cm}
{$x = 5\sin(10\pi t + \frac{\pi}{4})$ cm}
\loigiai{
Từ đồ thị: biên độ $A = 5$ cm.
 Vật đi từ biên dương (li độ $5\,\text{cm}$) đến li độ 0 mất $0,05\,\text{s}$, tức là $\frac{T}{4} = 0,05$ s, suy ra $T = 0,2$ s.
 Tần số góc:
\begin{aligned}
$\omega$ &= $\frac{2\pi}{T}$ = $\frac{2\pi}{0,2}$ = 10 $\pi\text{rad/s}$.
\end{aligned}
 Tại $t = 0$, li độ bằng biên độ dương, nên $\varphi_0 = 0$.
 Phương trình là $x = 5\cos(10\pi t)$ cm.
}
\end{ex}

% ===== Câu 195 =====
% ID: L11C1B1-59-TN
% Mức: VD
\begin{ex}
% ID: L11C1B1-59-TN
% Mức: VD
Đồ thị li độ - thời gian của một dao động điều hòa cho thấy vật đạt li độ cực đại $5\,\text{cm}$ tại $t = 0$, đạt li độ $0$ cm lần đầu tiên tại $t = 0,05$ s. Phương trình li độ là bao nhiêu?
\choice
{$x = 5\sin(10\pi t)$ cm}
{$x = 5\cos(10\pi t + \frac{\pi}{2})$ cm}
{\True $x = 5\cos(10\pi t)$ cm}
{$x = 5\sin(10\pi t + \frac{\pi}{4})$ cm}
\loigiai{
Từ đồ thị: biên độ $A = 5$ cm.
 Vật đi từ biên dương (li độ $5\,\text{cm}$) đến li độ 0 mất $0,05\,\text{s}$, tức là $\frac{T}{4} = 0,05$ s, suy ra $T = 0,2$ s.
 Tần số góc:
  \begin{aligned}
 $\omega$ &= $\frac{2\pi}{T}$ &= $\frac{2\pi}{0,2}$ &= 10 $\pi\\text{rad/s}$.
 \end{aligned} 
 Tại $t = 0$, li độ bằng biên độ dương, nên $\varphi_0 = 0$.
 Phương trình là $x = 5\cos(10\pi t)$ cm.
}
\end{ex}

% ===== Câu 196 =====
% ID: L11C1B1-60-TN
% Mức: VD
\begin{ex}
% ID: L11C1B1-60-TN
% Mức: VD
Đồ thị li độ - thời gian của một dao động điều hòa cho thấy vật đạt li độ cực đại $5\,\text{cm}$ tại $t = 0$, đạt li độ $0$ cm lần đầu tiên tại $t = 0,05$ s. Phương trình li độ là bao nhiêu?
\choice
{$x = 5\sin(10\pi t)$ cm}
{$x = 5\cos(10\pi t + \frac{\pi}{2})$ cm}
{\True $x = 5\cos(10\pi t)$ cm}
{$x = 5\sin(10\pi t + \frac{\pi}{4})$ cm}
\loigiai{
Từ đồ thị: biên độ $A = 5$ cm.
 Vật đi từ biên dương (li độ $5\,\text{cm}$) đến li độ 0 mất $0,05\,\text{s}$, tức là $\frac{T}{4} = 0,05$ s, suy ra $T = 0,2$ s.
 Tần số góc:
  \begin{aligned}
 $\omega$ &= $\frac{2\pi}{T}$ &= $\frac{2\pi}{0,2}$ &= 10 $\pi\\text{rad/s}$.
 \end{aligned} 
 Tại $t = 0$, li độ bằng biên độ dương, nên $\varphi_0 = 0$.
 Phương trình là $x = 5\cos(10\pi t)$ cm.
}
\end{ex}

% ===== Câu 197 =====
% ID: L11C1B1-61-TN
% Mức: VD
\begin{ex}
% ID: L11C1B1-61-TN
% Mức: VD
Đồ thị li độ-thời gian của một vật dao động điều hòa cho thấy: tại $t=0$, $x=4$ cm; tại $t=0,5$ s, $x=0$; tại $t=1$ s, $x=-4$ cm. Biên độ và chu kỳ của dao động là bao nhiêu? \nguon{Kết nối tri thức với cuộc sống - Vật lí 11, Bài 2}
\choice
{Biên độ $4\,\text{cm}$, chu kỳ $2\,\text{s}$}
{Biên độ $8\,\text{cm}$, chu kỳ $1\,\text{s}$}
{\True Biên độ $4\,\text{cm}$, chu kỳ $2\,\text{s}$}
{Biên độ $4\,\text{cm}$, chu kỳ $4\,\text{s}$}
\loigiai{
Từ đồ thị: tại $t=0$, $x=4$ cm và đây là giá trị dương lớn nhất trước khi dao động về 0 tại $t=0,5$ s, suy ra biên độ $A=4$ cm.
 Khoảng thời gian từ $t=0$ (vị trí $x=A$) đến $t=1$ s (vị trí $x=-A$) tương ứng với nửa chu kỳ.
  \begin{aligned}
 $\dfrac{T}{2}$ &= 1 $\\text{s}T$ &= 2 $\\text{s}$.
 \end{aligned} 
 Vậy biên độ là $4\,\text{cm}$, chu kỳ là $2\,\text{s}$.
 Chọn đáp án tương ứng.
}
\end{ex}

% ===== Câu 198 =====
% ID: L11C1B1-62-TN
% Mức: VD
\begin{ex}
% ID: L11C1B1-62-TN
% Mức: VD
Đồ thị li độ-thời gian của một vật dao động điều hòa cho thấy: tại $t=0$, $x=4$ cm; tại $t=0,5$ s, $x=0$; tại $t=1$ s, $x=-4$ cm. Biên độ và chu kỳ của dao động là bao nhiêu? \nguon{Kết nối tri thức với cuộc sống - Vật lí 11, Bài 2}
\choice
{Biên độ $4\,\text{cm}$, chu kỳ $2\,\text{s}$}
{Biên độ $8\,\text{cm}$, chu kỳ $1\,\text{s}$}
{\True Biên độ $4\,\text{cm}$, chu kỳ $2\,\text{s}$}
{Biên độ $4\,\text{cm}$, chu kỳ $4\,\text{s}$}
\loigiai{
Từ đồ thị: tại $t=0$, $x=4$ cm và đây là giá trị dương lớn nhất trước khi dao động về 0 tại $t=0,5$ s, suy ra biên độ $A=4$ cm.
 Khoảng thời gian từ $t=0$ (vị trí $x=A$) đến $t=1$ s (vị trí $x=-A$) tương ứng với nửa chu kỳ.
  \begin{aligned}
 $\dfrac{T}{2}$ &= 1 $\\text{s}T$ &= 2 $\\text{s}$.
 \end{aligned} 
 Vậy biên độ là $4\,\text{cm}$, chu kỳ là $2\,\text{s}$.
 Chọn đáp án tương ứng.
}
\end{ex}

% ===== Câu 199 =====
% ID: L11C1B1-63-TN
% Mức: VD
\begin{ex}
% ID: L11C1B1-63-TN
% Mức: VD
Đồ thị li độ-thời gian của một vật dao động điều hòa cho thấy: tại $t=0$, $x=4$ cm; tại $t=0,5$ s, $x=0$; tại $t=1$ s, $x=-4$ cm. Biên độ và chu kỳ của dao động là bao nhiêu? \nguon{Kết nối tri thức với cuộc sống - Vật lí 11, Bài 2}
\choice
{Biên độ $4\,\text{cm}$, chu kỳ $2\,\text{s}$}
{Biên độ $8\,\text{cm}$, chu kỳ $1\,\text{s}$}
{\True Biên độ $4\,\text{cm}$, chu kỳ $2\,\text{s}$}
{Biên độ $4\,\text{cm}$, chu kỳ $4\,\text{s}$}
\loigiai{
Từ đồ thị: tại $t=0$, $x=4$ cm và đây là giá trị dương lớn nhất trước khi dao động về 0 tại $t=0,5$ s, suy ra biên độ $A=4$ cm.
 Khoảng thời gian từ $t=0$ (vị trí $x=A$) đến $t=1$ s (vị trí $x=-A$) tương ứng với nửa chu kỳ.
  \begin{aligned}
 $\dfrac{T}{2}$ &= 1 $\\text{s}T$ &= 2 $\\text{s}$.
 \end{aligned} 
 Vậy biên độ là $4\,\text{cm}$, chu kỳ là $2\,\text{s}$.
 Chọn đáp án tương ứng.
}
\end{ex}
